\documentclass[12pt,a4paper]{article}

\usepackage[utf8]{inputenc}
\usepackage[T1]{fontenc}
\usepackage[spanish,es-nodecimaldot]{babel}
\usepackage{lmodern}
\usepackage{geometry}
\usepackage{graphicx}
\usepackage{array}
\usepackage{booktabs}
\usepackage{longtable}
\usepackage{tabularx}
\usepackage{enumitem}
\usepackage{tikz}
\usetikzlibrary{positioning}
\usepackage{xcolor}
\usepackage{hyperref}
\usepackage{fancyhdr}
\usepackage{float}

% Prefer ignored local data; fall back to the tracked non-sensitive example.
\IfFileExists{config/datos.local.tex}{\input{config/datos.local.tex}}{% Tracked, non-sensitive data template for local document customization.
% Copy this file to config/datos.local.tex and replace the placeholders.
% Active entry points load datos.local.tex when it exists; otherwise they
% load this example so the source remains structurally complete.

\newcommand{\PracticanteNombre}{Nombre del practicante}
\newcommand{\PracticanteCodigo}{CÓDIGO}
\newcommand{\PracticanteEspecialidad}{Ingeniería de Software}
\newcommand{\PracticanteFacultad}{FIMEES}
\newcommand{\PracticantePromocion}{AÑO}
\newcommand{\InstitucionNombre}{Dirección de Proyección Social y Extensión Cultural}
\newcommand{\InstitucionOficina}{Oficina de Dirección de Proyección Social y Extensión Cultural}
\newcommand{\InstitucionJefa}{Nombre de la jefa de la unidad receptora}
\newcommand{\DocenteNombre}{Nombre del docente de prácticas}
\newcommand{\FechaInicio}{DD/MM/AAAA}
\newcommand{\FechaFin}{DD/MM/AAAA}
\newcommand{\CentroDireccion}{Dirección del centro de prácticas}
\newcommand{\TelefonoMelissa}{Teléfono de contacto del proyecto}
\newcommand{\TelefonoJefa}{Teléfono de la dirección}
\newcommand{\TelefonoOficina}{Teléfono de la oficina}
\newcommand{\SitioWeb}{https://example.invalid}
}

\geometry{left=2.5cm,right=2.5cm,top=2.5cm,bottom=2.5cm}
\hypersetup{colorlinks=true,linkcolor=black,urlcolor=blue}
\setlength{\parindent}{1.25cm}
\setlength{\parskip}{0.35em}
\renewcommand{\arraystretch}{1.35}
\newcolumntype{Y}{>{\raggedright\arraybackslash}X}
\newcommand{\checkmarkx}{\textbf{X}}

\pagestyle{fancy}
\fancyhf{}
\lhead{Plan de Prácticas Pre Profesionales II}
\rhead{\thepage}
\renewcommand{\headrulewidth}{0.4pt}

\begin{document}

% --------------------------------------------------------------------------
% PORTADA
% --------------------------------------------------------------------------
\begin{titlepage}
\centering
\thispagestyle{empty}

{\Large\textbf{UNIVERSIDAD NACIONAL DEL ALTIPLANO}}\\[0.35cm]
{\large\textbf{FACULTAD DE INGENIERÍA ELÉCTRICA, ELECTRÓNICA Y SISTEMAS}}\\[0.35cm]
{\large\textbf{ESCUELA PROFESIONAL DE INGENIERÍA DE SISTEMAS}}\\[1.4cm]

\begin{tikzpicture}
  \draw[fill=green!55!black,draw=green!55!black] (0,0) circle (1.7cm);
  \node[align=center,text=white,font=\bfseries\small] at (0,0)
    {UNA\\PUNO};
\end{tikzpicture}\\[1.2cm]

{\Large\textbf{PLAN DE PRÁCTICAS PRE PROFESIONALES II}}\\[1.5cm]
{\large\textbf{CENTRO DE PRÁCTICAS:}}\\[0.2cm]
{\large \InstitucionNombre}\\[0.9cm]
{\large\textbf{PRESENTADO POR:}}\\[0.2cm]
{\large \PracticanteNombre}\\[0.9cm]
{\large\textbf{DOCENTE:}}\\[0.2cm]
{\large \DocenteNombre}\\[0.9cm]
{\large\textbf{SEMESTRE:} X \hspace{1.5cm} \textbf{GRUPO:} B}\\[1.2cm]

{\large PUNO -- PERÚ}\\[1.5cm]

\textbf{FIRMAS DE CONFORMIDAD DEL PLAN DE PRÁCTICAS}\
\textbf{PRE PROFESIONALES II}\\[1cm]

\begin{tabular}{>{\centering\arraybackslash}p{0.4\textwidth}>{\centering\arraybackslash}p{0.4\textwidth}}
\rule{5cm}{0.4pt} & \rule{5cm}{0.4pt}\\
\PracticanteNombre & \InstitucionJefa\\
Practicante preprofesional & Jefa de la unidad receptora
\end{tabular}

\vfill
\textbf{Docente de prácticas preprofesionales}\\[0.7cm]
\rule{6cm}{0.4pt}\\
\DocenteNombre
\end{titlepage}

\tableofcontents
\newpage

% --------------------------------------------------------------------------
\section{Parte informativa}

\begin{description}[style=multiline,leftmargin=4.5cm,labelwidth=4cm]
  \item[1.1. Nombre del practicante:] \PracticanteNombre.
  \item[1.2. Especialidad de prácticas:] \PracticanteEspecialidad.
  \item[1.3. Nombre de la institución:] \InstitucionNombre.
  \item[1.4. Directora de la institución:] \InstitucionJefa.
  \item[1.5. Jefa de la unidad receptora:] \InstitucionJefa.
  \item[1.6. Duración:] Cuatro meses.
  \item[Inicio:] \FechaInicio.
  \item[Finalización:] \FechaFin.
  \item[1.7. Docente de prácticas:] \DocenteNombre.
  \item[1.8. Programa de estudios:] Ingeniería de Sistemas.
  \item[1.9. Facultad:] Ingeniería Eléctrica, Electrónica y Sistemas (FIMEES).
\end{description}

\section{Del centro de prácticas}

\subsection{Denominación del centro de prácticas}
Oficina de Proyección Social y Extensión Cultural (DPSEC).

\subsection{Dirección del centro de prácticas}
Auditorio Mango UNA Puno, tercer piso.

\subsection{Teléfono, correo electrónico y página web}
\begin{itemize}
  \item \TelefonoMelissa: Ing. Melissa, encargada del proyecto.
  \item \TelefonoJefa: \InstitucionJefa, directora de la oficina DPSEC.
  \item \TelefonoOficina: Oficina DPSEC, secretaría.
  \item Página web: \url{\SitioWeb}
\end{itemize}

\subsection{Visión}
La Dirección de Proyección Social y Extensión Cultural fortalece el desarrollo de la región bajo los principios de la ciencia, la tecnología, los valores culturales y la investigación, en el marco de una responsabilidad social sostenible.

\subsection{Misión}
La Dirección de Proyección Social y Extensión Cultural genera interacción con la comunidad para el desarrollo regional a través de la ciencia, la tecnología y las expresiones culturales como parte de la labor académica y de investigación, en el marco de una proyección social y extensión cultural sostenible.

\subsection{Valores institucionales}
El centro de prácticas promueve la responsabilidad social, la colaboración, la ética profesional, la innovación, el compromiso con la comunidad y el respeto por la diversidad cultural.

\subsection{Objetivos de la institución}
Contribuir al desarrollo regional mediante la articulación de la universidad con la comunidad, impulsando proyectos sociales, culturales, tecnológicos y ambientales.

\subsection{Actividad principal de la institución}
Promover la interacción entre la universidad y la comunidad a través de proyectos sociales, culturales y de extensión, contribuyendo al desarrollo regional.

\subsection{Organigrama}
La estructura funcional del centro de prácticas se puede representar de la siguiente manera:

\begin{center}
\begin{tikzpicture}[node distance=7mm and 8mm, every node/.style={draw, rounded corners, align=center, font=\small, minimum height=8mm}]
  \node[fill=green!20] (vracad) {VRACAD -- UNA};
  \node[fill=green!45,below=of vracad,minimum width=7cm] (dpsec) {Dirección de Proyección Social\\Extensión Cultural};
  \node[below left=of dpsec,xshift=-15mm] (admin) {Secretaría\\administrativa};
  \node[below right=of dpsec,xshift=15mm] (com) {Área de\\comunicación};
  \node[fill=green!20,below=18mm of dpsec,xshift=-5cm] (social) {Subunidad de\\proyección social\\y extensión cultural};
  \node[fill=green!20,below=18mm of dpsec] (graduados) {Subunidad de seguimiento\\de graduados};
  \node[fill=green!20,below=18mm of dpsec,xshift=5cm] (ambiental) {Subunidad de\\gestión ambiental};
  \draw (vracad)--(dpsec);
  \draw (dpsec)--(admin) (dpsec)--(com);
  \draw (dpsec)--(social) (dpsec)--(graduados) (dpsec)--(ambiental);
\end{tikzpicture}
\end{center}

\subsection{Funciones de la oficina, área o dependencia de prácticas}
La DPSEC coordina y acompaña proyectos de proyección social, extensión cultural, seguimiento de graduados y gestión ambiental, articulando las actividades universitarias con las necesidades de la comunidad.

\newpage
\section{Objetivos generales y específicos de las prácticas}

\subsection{Objetivo general}
Mejorar el sistema web Genseg para la gestión de proyectos y actividades de la Oficina de Proyección Social y Extensión Cultural (DPSEC), aplicando conocimientos de ingeniería de software.

\subsection{Objetivos específicos}
\begin{itemize}
  \item Analizar los requerimientos funcionales y no funcionales del sistema Genseg versión 1 en colaboración con el personal de la DPSEC.
  \item Mejorar la arquitectura del sistema, incluyendo la base de datos y los componentes de backend y frontend.
  \item Desarrollar el backend utilizando Node.js, TypeScript y Prisma, implementando la lógica de negocio y la API REST.
  \item Desarrollar el frontend utilizando Next.js y TypeScript, siguiendo el diseño propuesto en Figma para asegurar una interfaz intuitiva y funcional.
  \item Realizar pruebas unitarias y de integración para asegurar la calidad y estabilidad del sistema.
  \item Documentar el proceso de desarrollo y preparar el sistema para su despliegue.
\end{itemize}

\section{Componentes del plan}

\subsection{Función principal del puesto}
Practicante de Desarrollo de Software.

\subsection{Actividades y tareas principales}
\small
\begin{longtable}{p{0.22\textwidth}p{0.15\textwidth}p{0.14\textwidth}p{0.27\textwidth}p{0.14\textwidth}}
\caption{Actividades principales del plan}\label{tab:actividades}\\
\toprule
\textbf{Actividad} & \textbf{Unidad de medida} & \textbf{Meta} & \textbf{Herramienta} & \textbf{Responsable}\\
\midrule
\endfirsthead
\toprule
\textbf{Actividad} & \textbf{Unidad de medida} & \textbf{Meta} & \textbf{Herramienta} & \textbf{Responsable}\\
\midrule
\endhead
Análisis de requerimientos y diseño del sistema & Documento & En curso & Figma y documentación & \PracticanteNombre\\
Desarrollo del backend (API REST y lógica) & Módulos & Pendiente & Node.js, TypeScript, Prisma, PostgreSQL y Docker & \PracticanteNombre\\
Desarrollo del frontend (interfaz de usuario) & Vistas y componentes & Pendiente & Next.js, TypeScript, Tailwind CSS y Figma & \PracticanteNombre\\
Pruebas unitarias y de integración & Casos de prueba & Pendiente & Herramientas de testing y Postman & \PracticanteNombre\\
Documentación del sistema & Manual e informe & Pendiente & Word & \PracticanteNombre\\
\bottomrule
\end{longtable}
\normalsize

\subsection{Cronograma}
\small
\begin{tabularx}{\textwidth}{p{0.27\textwidth}*{16}{>{\centering\arraybackslash}X}}
\toprule
\textbf{Tarea} & 1&2&3&4&5&6&7&8&9&10&11&12&13&14&15&16\\
\midrule
Análisis de requerimientos & \checkmarkx&\checkmarkx&\checkmarkx&\checkmarkx&&&&&&&&&&&&\\
Diseño & &&&&\checkmarkx&\checkmarkx&&&&&&&&&&&&\\
Desarrollo backend & &&&&&&\checkmarkx&\checkmarkx&\checkmarkx&\checkmarkx&&&&&&&&\\
Desarrollo frontend & &&&&&&&&&&\checkmarkx&\checkmarkx&\checkmarkx&\checkmarkx&&&\\
Integración backend--frontend & &&&&&&&&&&&&&&\checkmarkx&\checkmarkx\\
Pruebas y ajustes & &&&&&&&&&&&&&&\checkmarkx&\checkmarkx\\
Documentación y preparación de entrega & &&&&&&&&&&&&&&\checkmarkx&\checkmarkx\\
\bottomrule
\end{tabularx}
\normalsize

\subsection{Competencias}
\begin{tabularx}{\textwidth}{>{\bfseries}Y Y}
\toprule
Competencia general & Indicador de logro\\
\midrule
Análisis y solución de problemas & Identifica necesidades, analiza requerimientos y propone soluciones de software.\\
Trabajo en equipo y comunicación & Participa en reuniones, documenta acuerdos y coordina con el personal de la DPSEC.\\
Desarrollo tecnológico & Implementa componentes backend y frontend aplicando buenas prácticas de ingeniería de software.\\
Responsabilidad profesional & Cumple el cronograma, documenta los avances y respeta las normas de la institución.\\
\bottomrule
\end{tabularx}

\section{Supervisión}
La supervisión está a cargo de \InstitucionJefa y del docente del curso de prácticas, \DocenteNombre.

\section{Horario de prácticas}
\begin{center}
\begin{tabular}{>{\bfseries}p{3cm}*{5}{>{\centering\arraybackslash}p{2cm}}}
\toprule
Horario & Lunes & Martes & Miércoles & Jueves & Viernes\\
\midrule
8:00--12:00 & DPSEC & DPSEC & DPSEC & DPSEC & DPSEC\\
\bottomrule
\end{tabular}
\end{center}

\section{Calificación y requisitos de aprobación}

\subsection{Criterios de calificación}
\begin{enumerate}
  \item \textbf{Asistencia:} se evaluará la puntualidad y el cumplimiento de las horas de prácticas. Se requiere un mínimo de 80\% de asistencia.
  \item \textbf{Cumplimiento de actividades:} el practicante debe cumplir las actividades y tareas asignadas, demostrando habilidades técnicas y actitud proactiva.
  \item \textbf{Desempeño profesional:} se evaluará la actitud, responsabilidad y comportamiento ético durante el desarrollo de las tareas.
  \item \textbf{Informe final:} se presentará un informe que documente las actividades realizadas, el aprendizaje y la experiencia adquirida.
  \item \textbf{Evaluación del supervisor:} el supervisor asignará una calificación basada en la ejecución de las tareas, las competencias demostradas y el impacto de la contribución del practicante.
\end{enumerate}

\subsection{Requisitos para la aprobación}
\begin{enumerate}
  \item Cumplir una asistencia mínima del 80\% a las prácticas preprofesionales.
  \item Cumplir las metas establecidas en el plan de trabajo, incluida la presentación y aprobación del informe final.
  \item Obtener una evaluación positiva por parte del supervisor de la institución.
  \item Presentar y aprobar el informe de prácticas dentro del plazo establecido por la institución educativa.
\end{enumerate}

\section{Reglamento e instrumentos generales}

\subsection*{Artículo 5: objetivos de las prácticas preprofesionales}
\begin{enumerate}[label=\Alph*.]
  \item Aplicar los conocimientos teóricos y prácticos adquiridos durante la carrera a la realidad profesional.
  \item Cumplir uno de los requisitos establecidos en la currícula para obtener la condición de egresado.
  \item Afianzar conocimientos y experiencias, colaborando con la formación integral del practicante.
  \item Adquirir habilidades y destrezas en la planificación, desarrollo y evaluación de programas propios de la escuela profesional.
  \item Aprender a trabajar en distintos niveles como miembro de un equipo interdisciplinario y multifacético.
  \item Realizar actividades de extensión universitaria en beneficio de la comunidad y de la región.
\end{enumerate}

\subsection*{Artículo 19: causas de sanción}
\begin{enumerate}
  \item Demora injustificada en integrarse al centro de prácticas.
  \item No ejecutar las prácticas por motivos injustificados.
  \item Abandonar las prácticas.
  \item Presentar el informe de prácticas fuera del plazo establecido.
  \item Desconocer o no respetar las normas organizacionales o laborales del centro de prácticas.
\end{enumerate}

\subsection*{Artículo 20}
Las sanciones por las faltas ocasionan, en primer lugar, una amonestación; luego, la separación de las prácticas y, finalmente, la anulación de estas, previa investigación de las causas.

\subsection*{Artículo 21}
Los estudiantes sancionados con la anulación de la práctica preprofesional tendrán una última oportunidad de repetirla después de una convocatoria.

\section{Referencias}
\begin{itemize}
  \item Dirección de Proyección Social y Extensión Cultural. Información institucional y documentación del centro de prácticas.
  \item Universidad Nacional del Altiplano. Reglamento e instrumentos generales de prácticas preprofesionales.
\end{itemize}

\end{document}
